\documentclass[11pt]{article}
\usepackage{siunitx}
\usepackage{cleveref}
\usepackage[margin=1in]{geometry}
\usepackage{graphicx}
\usepackage{booktabs}
\usepackage{float}
\usepackage{caption}
\usepackage{subcaption}

% Change only these lines to point the note at another rendered report.
% Change only these lines to point the note at another rendered report.
\newcommand{\reportroot}{../Compute-Link-Attenuations/HundredPatches/norm/report}
\newcommand{\datasetlabel}{HundredPatches/norm}
\newcommand{\distanceboxplot}{\reportroot/images/distance_profiles_box_whisker/linear/distance_iqr_medians_rainy_multi_k3_box_whisker.png}
\newcommand{\nonrainydistanceboxplot}{\reportroot/images/distance_profiles_box_whisker/linear/distance_iqr_medians_nonrainy_multi_k3_box_whisker.png}
\newcommand{\idwconfusionplot}{\reportroot/images/confusion_maps/IDW_confusion_map.png}
\newcommand{\ildwconfusionplot}{\reportroot/images/confusion_maps/ILDW_confusion_map.png}
\newcommand{\nonlinearconfusionplot}{\reportroot/images/confusion_maps/Nonlinear_Optimizer_confusion_map.png}
\newcommand{\patchgeometryplot}{\reportroot/images/patch_overview/patch_dimensions_scatter.png}
\newcommand{\linkhistplot}{\reportroot/images/patch_overview/link_count_histogram.png}
\newcommand{\largestpatchbinsplot}{\reportroot/images/largest_patch_distance_bins/largest_patch_distance_bins_k3.png}
\newcommand{\contextroot}{\reportroot/images/rain_patch_explorer_maps}
\newcommand{\contextscreenshotone}{\contextroot/202301280900_patch000_map.png}
\newcommand{\contextscreenshottwo}{\contextroot/202301310900_patch000_map.png}
\newcommand{\contextscreenshotthree}{\contextroot/202301311600_patch000_map.png}
\newcommand{\idwcaseone}{\reportroot/images/patch_error_maps/IDW/RAD_OPERA_HOURLY_RAINFALL_ACCUMULATION_202301280900_patch000_rainy.png}
\newcommand{\ildwcaseone}{\reportroot/images/patch_error_maps/ILDW/RAD_OPERA_HOURLY_RAINFALL_ACCUMULATION_202301280900_patch000_rainy.png}
\newcommand{\optcaseone}{\reportroot/images/patch_error_maps/OPT_NORM_ILDW_MULT_ILDW_INIT/RAD_OPERA_HOURLY_RAINFALL_ACCUMULATION_202301280900_patch000_rainy.png}
\newcommand{\idwcasetwo}{\reportroot/images/patch_error_maps/IDW/RAD_OPERA_HOURLY_RAINFALL_ACCUMULATION_202301310900_patch000_rainy.png}
\newcommand{\ildwcasetwo}{\reportroot/images/patch_error_maps/ILDW/RAD_OPERA_HOURLY_RAINFALL_ACCUMULATION_202301310900_patch000_rainy.png}
\newcommand{\optcasetwo}{\reportroot/images/patch_error_maps/OPT_NORM_ILDW_MULT_ILDW_INIT/RAD_OPERA_HOURLY_RAINFALL_ACCUMULATION_202301310900_patch000_rainy.png}
\newcommand{\idwcasethree}{\reportroot/images/patch_error_maps/IDW/RAD_OPERA_HOURLY_RAINFALL_ACCUMULATION_202301311600_patch000_rainy.png}
\newcommand{\ildwcasethree}{\reportroot/images/patch_error_maps/ILDW/RAD_OPERA_HOURLY_RAINFALL_ACCUMULATION_202301311600_patch000_rainy.png}
\newcommand{\optcasethree}{\reportroot/images/patch_error_maps/OPT_NORM_ILDW_MULT_ILDW_INIT/RAD_OPERA_HOURLY_RAINFALL_ACCUMULATION_202301311600_patch000_rainy.png}


\title{Rainfall Estimation}
\author{}
\date{}

\begin{document}
\maketitle

This note summarizes results for a collection of 100 patches. All aggregate figures and tables that follow should therefore be read as patch-level summaries over this hundred-patch evaluation set, rather than as statistics from a single patch or from the full archive of available patches.

\section*{Relative Absolute Error Rainy-Pixel Distance Profile}
\begin{figure}[H]
  \centering
  \includegraphics[width=\textwidth]{\distanceboxplot}
  \caption{Rainy-pixel distance-profile box-and-whisker summary for distance to the third closest link. Each patch contributes one median rainy-pixel relative-absolute-error (RAE) value per distance bin, and these patch-level medians are then summarized across patches with a box-and-whisker plot. Here, a rainy pixel is a pixel whose ground-truth rainfall is at least the configured wet threshold (set to 0.6 mm/h), and relative absolute error for a rainy pixel $p$ is defined as $|GT(p)-PRED(p)| / GT(p)$.}
  \label{fig:RAE}
\end{figure}

This plot summarizes how solver performance changes as pixels get farther from the third-nearest link. Lower values indicate lower rainy-pixel relative absolute error in that distance bin.

\section*{Absolute Error Non-Rainy-Pixel Distance Profile}
\begin{figure}[H]
  \centering
  \includegraphics[width=\textwidth]{\nonrainydistanceboxplot }
  \caption{Non-rainy-pixel distance-profile box-and-whisker summary for distance to the third closest link. Each patch contributes one median non-rainy-pixel absolute error value per distance bin, and these patch-level medians are then summarized across patches with a box-and-whisker plot. Here, a non-rainy pixel is a pixel whose ground-truth rainfall is below the configured wet threshold (set to 0.6 mm/h), and the absolute error for a non-rainy pixel $p$ is defined as $|GT(p)-PRED(p)|$.}
  \label{fig:AE}
\end{figure}


\subsection*{Rainy-pixel count table for distance to the third closest link}
The table below gives the average number of rainy pixels per patch and the corresponding standard deviation for each distance bin to the 3rd closest link. These are the same patch-level counts summarized in the tick labels of the rainy-pixel distance-profile plots.

\begin{table}[H]
  \centering
    \begin{tabular}{l
                  l
                  S[table-format=5.1]}
    \toprule
    {Distance to third closest link (m)} & {Mean count (\% of total)} & {SD} \\
    \midrule
    $[0,125]$       & 4267.0 (1.79\%) &   960.4 \\
    $(125,375]$     & 16737.7 (7.01\%) &  3803.2 \\
    $(375,750]$     & 26809.9 (11.22\%) &  6707.0 \\
    $(750,1500]$    & 51988.5 (21.76\%) & 14468.0 \\
    $(1500,3125]$   & 88839.4 (37.18\%) & 27531.1 \\
    $(3125,6000]$   & 45929.4 (19.22\%) & 15679.4 \\
    $(6000,9000]$   & 3835.8 (1.61\%) &  2725.7 \\
    $(9000,\infty)$ & 514.4 (0.22\%) &   966.9 \\
    \midrule
    {Total}         & 238922.1 (100\%) & {} \\
    \bottomrule
  \end{tabular}
  \caption{Distribution of the number of pixels per distance bin. In each patch, pixels are classified into bins according to their distance to the 3rd-closest link. Formally, a pixel $p$ belongs to bin $(d_1,d_2]$ if the disc of radius $d_2$ centered at $p$ intersects at least $3$ links and the disc of radius $d_1$ centered at $p$ intersects at most $2$ links. Let $|\mathrm{bin}_i(P)|$ denote the number of pixels in patch $P$ that belong to $\mathrm{bin}_i$. The table reports, for each bin, the mean and standard deviation of $|\mathrm{bin}_i(P)|$ across patches. The last bin, 
$(9000,\infty)$ is merged with the previous bin in \Cref{fig:RAE} because it is  sparsely populated.}
  %\caption{Rainy-pixel counts per patch for the distance bins to the third closest link used in the rainy-pixel distance-profile plots.}
\end{table}

\section*{Confusion Matrix}
\begin{table}[H]
  \centering
  \begin{minipage}[t]{0.31\textwidth}
    \centering
    \textbf{IDW}

    \vspace{0.35em}
    \scriptsize
    \begin{tabular}{lcc}
      \toprule
       & \multicolumn{2}{c}{Prediction} \\
      GT & Wet & Dry \\
      \midrule
      Wet & $0.893 \pm 0.013$ & $0.002 \pm 0.000$ \\
      Dry & $0.040 \pm 0.004$ & $0.065 \pm 0.010$ \\
      \bottomrule
    \end{tabular}
  \end{minipage}
  \hfill
  \begin{minipage}[t]{0.31\textwidth}
    \centering
    \textbf{ILDW}

    \vspace{0.35em}
    \scriptsize
    \begin{tabular}{lcc}
      \toprule
       & \multicolumn{2}{c}{Prediction} \\
      GT & Wet & Dry \\
      \midrule
      Wet & $0.892 \pm 0.013$ & $0.003 \pm 0.000$ \\
      Dry & $0.037 \pm 0.004$ & $0.068 \pm 0.010$ \\
      \bottomrule
    \end{tabular}
  \end{minipage}
  \hfill
  \begin{minipage}[t]{0.31\textwidth}
    \centering
    \textbf{Nonlinear Optimizer}

    \vspace{0.35em}
    \scriptsize
    \begin{tabular}{lcc}
      \toprule
       & \multicolumn{2}{c}{Prediction} \\
      GT & Wet & Dry \\
      \midrule
      Wet & $0.890 \pm 0.013$ & $0.005 \pm 0.001$ \\
      Dry & $0.030 \pm 0.004$ & $0.075 \pm 0.010$ \\
      \bottomrule
    \end{tabular}
  \end{minipage}
  \caption{Patch-averaged wet/dry confusion matrices for IDW, ILDW, and the nonlinear optimizer at a threshold of $0.6$ mm/h, where a pixel is called wet if its rainfall is at least $0.6$ mm/h and dry otherwise. For each patch, pixels are first assigned to the four confusion categories defined by the ground-truth and predicted wet/dry labels: top-left is TP, top-right is FN, bottom-left is FP, and bottom-right is TN. The count in each category is then divided by the total number of pixels in that patch, producing a patch-level pixel fraction. The table entries report the mean of those patch-level fractions across the 100 evaluated patches, together with the standard error of the mean (SEM).}
\end{table}


\section*{Patch Geometry Overview}
\begin{figure}[H]
  \centering
  \includegraphics[width=0.82\textwidth]{\patchgeometryplot}
  \caption{Patch width versus patch height for the 100 patches, with one dot per patch.  Across patches, the mean patch area is 269,491 pixels with SD of 76,755 pixels (pixels are squares with side length $125$m).}
 % \caption{Patch width versus patch height, with one marker per patch. The footnote on the figure reports the average patch size in pixels and its standard deviation across patches.}
\end{figure}

% This figure makes it easy to see whether the selected patches occupy a tight geometry range or whether the report is mixing very different patch extents.

% \section*{Link Count Distribution}
% \begin{figure}[H]
%   \centering
%   \includegraphics[width=0.82\textwidth]{\linkhistplot}
%   \caption{Histogram of the number of links per patch. The figure footnote also reports the average parallel-link ratio, defined here as the fraction of links in a patch that have at least one parallel partner.}
% \end{figure}

% This histogram summarizes how link density varies across the selected patches. The parallel-link ratio is included because the current expectation is that it should stay close to $1$ for this dataset.

\begin{table}[H]
  \centering
  \begin{tabular}{lrrrrr}
    \toprule
    Patch size (km) & \# links & 1-multiplicity & 2-multiplicity \\
    \midrule
    $50 \times 50$ & 939 & 99 & 420 \\
    $62 \times 60$ & 1254 & 128 & 563  \\
    $72 \times 70$ & 1507 & 157 & 675  \\
    \bottomrule
  \end{tabular}
  \caption{Representative link multiplicities for three patch dimensions.
The multiplicity of a link is defined as the number of frequencies assigned to it.
The second column reports the total number of link–frequency pairs.
In the data set, link multiplicities were either 1 or 2, with approximately $10\%$ of links having multiplicity 1.
}

\end{table}

\section*{$J_{\mathrm{atten}}$ Summary}
\begin{table}[H]
  \centering
  \begin{tabular}{lcc}
    \toprule
    Solver & Mean $J_{\mathrm{atten}}$ (dB$^2$/km) & SD \\
    \midrule
    IDW & 0.01522 & 0.03808 \\
    ILDW & 0.00289 & 0.00757 \\
    Nonlinear Optimizer & 0.00020 & 0.00103 \\
    \bottomrule
  \end{tabular}
  \caption{Patch-level summary of $J_{\mathrm{atten}}$ across the 100 evaluated patches. For a given patch, $J_{\mathrm{atten}}$ is the attenuation-fit term $J_{\mathrm{atten}} = \frac{1}{N_{\mathrm{links}}}\sum_{\ell \in \mathrm{links}} \frac{\left(\hat{A}_\ell - A_\ell^{\mathrm{obs}}\right)^2}{|\ell|}$, where $\hat{A}_\ell$ is the attenuation implied by the estimated rainfall field on link $\ell$, $A_\ell^{\mathrm{obs}}$ is the observed attenuation, and $|\ell|$ is the link length in kilometers. The table reports the mean and standard deviation of this per-patch quantity for each solver. Lower values indicate better agreement with the observed link attenuations.}
\end{table}

\section*{Absolute Attenuation Error per km}
\begin{table}[H]
  \centering
  \begin{tabular}{lcc}
    \toprule
    Solver & Mean (dB/km) & SD \\
    \midrule
    IDW & 0.05034 & 0.03413 \\
    ILDW & 0.01356 & 0.00948 \\
    Nonlinear Optimizer & 0.00230 & 0.00437 \\
    \bottomrule
  \end{tabular}
  \caption{Patch-level summary of the mean absolute attenuation error per km. For a given patch, we compute $\frac{1}{N_{\mathrm{links}}}\sum_{\ell \in \mathrm{links}} \frac{\left|\hat{A}_\ell - A_\ell^{\mathrm{obs}}\right|}{|\ell|}$, where $\hat{A}_\ell$ is the attenuation induced by the estimated rainfall field on link $\ell$, $A_\ell^{\mathrm{obs}}$ is the observed attenuation, and $|\ell|$ is the link length in kilometers. Thus, the absolute attenuation error of every link is normalized by the link's length. The table reports the mean and standard deviation of this quantity across the $100$ evaluated patches for each solver. Lower values indicate better agreement with the observed link attenuations.
  }
\end{table}

% \begin{table}[H]
%   \centering
%   \begin{tabular}{lcc}
%     \toprule
%     Solver & Mean & SD \\
%     \midrule
%     IDW & 0.05512 & 0.04126 \\
%     ILDW & 0.00877 & 0.00700 \\
%     Nonlinear Optimizer & 0.00211 & 0.00365 \\
%     \bottomrule
%   \end{tabular}
%   \caption{Patch-level summary of the globally length-normalized absolute attenuation error. For a given patch, we compute $\frac{\sum_{k \in \mathrm{links}} \left|\hat{A}_k - A_k^{\mathrm{obs}}\right|}{\sum_{k \in \mathrm{links}} L_k^{\mathrm{km}}}$ using the same notation as above. Unlike the previous table, this metric first sums the absolute attenuation error over all valid links in the patch and then divides by the total valid-link length, so longer links receive more weight. The table reports the mean and standard deviation of this per-patch quantity across the 100 evaluated patches for each solver. Lower values again indicate better agreement with the observed link attenuations.}
% \end{table}

\end{document}
